\section*{第 \ref{ch:g12:deriv} 章 --- 求导与凸性}

\begin{solution}{exo:g12:deriv:1}
$f'(x) = 3x^2 - 5$。

商法则：
$g'(x) = \dfrac{(x^2+1) - x\cdot 2x}{(x^2+1)^2} = \dfrac{1 - x^2}{(x^2+1)^2}$。

取 $u = 2x+1$，由链式法则 $h'(x) = 5 \cdot 2 \cdot (2x+1)^4 = 10(2x+1)^4$。

取 $u = x^2 + 1$，由链式法则：
$k'(x) = \dfrac{2x}{2\sqrt{x^2+1}} = \dfrac{x}{\sqrt{x^2+1}}$。
\end{solution}

\begin{solution}{exo:g12:deriv:2}
$f'(x) = 2x - 3$，故 $f(2) = -1$，$f'(2) = 1$，切线为
$y = -1 + (x - 2) = x - 3$。作差得
\[
f(x) - (x - 3) = x^2 - 4x + 4 = (x-2)^2 \geq 0 ,
\]
因此曲线处处位于切线上方，仅在 $x = 2$
处与切线相切（与预期一致：$f$ 为凸函数）。
\end{solution}

\begin{solution}{exo:g12:deriv:3}
两者都是差商。取 $f(x) = (1+x)^{100}$，考虑 $0$ 处：
$f'(x) = 100(1+x)^{99}$，故极限为 $f'(0) = 100$。
取 $g(x) = \sqrt{x}$ 并考虑 $4$ 处：$g'(x) = \frac{1}{2\sqrt{x}}$，故极限为
$g'(4) = \frac14$。
\end{solution}

\begin{solution}{exo:g12:deriv:4}
定义域为 $\R \setminus \{1\}$。极限：在 $\pm\infty$ 处，$f(x) \sim x \to
\pm\infty$；在 $1^\pm$ 处，分子 $\to 4 > 0$，分母 $\to 0^\pm$，
故 $f \to \pm\infty$；直线 $x = 1$ 为铅直渐近线。

\[
f'(x) = \frac{2x(x-1) - (x^2+3)}{(x-1)^2} = \frac{x^2 - 2x - 3}{(x-1)^2}
= \frac{(x-3)(x+1)}{(x-1)^2}.
\]
于是 $f' > 0$ 在 $\intoo{-\infty}{-1}$ 与 $\intoo{3}{+\infty}$ 上，$f' < 0$ 在
$\intoo{-1}{1}$ 与 $\intoo{1}{3}$ 上。函数先递增，
至极大值 $f(-1) = -2$，递减至 $-\infty$，越过渐近线后跳至 $+\infty$，
再递减至极小值 $f(3) = 6$，然后递增。
\end{solution}

\begin{solution}{exo:g12:deriv:5}
当 $x \in \intoo{0}{15}$ 时，盒子底面为边长 $30 - 2x$ 的正方形，
高为 $x$，故
\[
V(x) = x(30 - 2x)^2 .
\]
$V'(x) = (30-2x)^2 + x \cdot 2(30-2x)(-2) = (30-2x)(30 - 2x - 4x)
= (30-2x)(30-6x)$。在 $\intoo{0}{15}$ 上 $30 - 2x > 0$，故 $V'$ 的符号与
$30 - 6x$ 相同：在 $x = 5$ 之前为正，之后为负。体积在
$x = 5$ 厘米时最大，此时 $V(5) = 5 \times 20^2 = 2000\ \text{cm}^3$。
\end{solution}

\begin{solution}{exo:g12:deriv:6}
\emph{1.} $f'(x) = 4x^3 - 12x^2$，$f''(x) = 12x^2 - 24x = 12x(x-2)$。
于是 $f'' > 0$ 在 $\intoo{-\infty}{0}$ 与 $\intoo{2}{+\infty}$ 上，故曲线为凸，
$f'' < 0$ 在 $\intoo{0}{2}$ 上，曲线为凹。$f''$ 在 $0$ 与 $2$ 处变号，
故二者均为拐点。

\emph{2.} 取 $a = 2$：$f(2) = 16 - 32 + 10 = -6$，$f'(2) = 32 - 48 = -16$，
切线为 $y = -6 - 16(x-2) = -16x + 26$。作差得
\[
g(x) = f(x) + 16x - 26 = x^4 - 4x^3 + 16x - 16 .
\]
由 $g(2) = g'(2) = g''(2) = 0$ 知 $(x-2)^3$ 整除 $g$，作除法得
$g(x) = (x-2)^3(x+2)$。在 $x = 2$ 附近 $x + 2 > 0$，故 $g$ 的符号与
$(x-2)^3$ 相同：$2$ 之前为负，之后为正。曲线穿过
切线，因此 $2$ 确实是拐点。
\end{solution}

\begin{solution}{exo:g12:deriv:7}
函数 $f(x) = \sqrt{x}$ 满足
$f''(x) = -\frac{1}{4}x^{-3/2} < 0$ 在 $\intoo{0}{+\infty}$ 上，故曲线为凹，
其图像位于每条切线的 \emph{下方}。在 $a = 1$ 处，切线为
\[
y = f(1) + f'(1)(x - 1) = 1 + \tfrac12 (x-1) = \frac{x+1}{2},
\]
于是 $\sqrt{x} \leq \frac{x+1}{2}$ 对所有 $x > 0$ 成立。等号成立时，
曲线与切线相切，这仅在切点
$x = 1$ 处发生。（等价地，$\left(\sqrt x - 1\right)^2 \geq 0$。）
\end{solution}

\begin{solution}{exo:g12:deriv:8}
由 \cref{thm:g12:deriv:convexity} 可知，可导凸函数的图像
位于其每条切线的上方。在 $a$ 处的切线水平
（$f'(a) = 0$），方程为 $y = f(a)$。于是 $f(x) \geq f(a)$ 对所有
$x \in \R$ 成立，即 $f(a)$ 为全局最小值。
\end{solution}

\begin{solution}{exo:g12:deriv:9}
\emph{周长固定。} 边长为 $x$ 与 $\frac{p}{2} - x$ 的矩形
（$0 < x < \frac p2$）面积为 $S(x) = x\left(\frac p2 - x\right)$。于是
$S'(x) = \frac p2 - 2x$ 在 $x = \frac p4$ 处为零，之前为正，
之后为负，故最大值在 $x = \frac{p}{4}$ 处取得，此时两边长均为
$\frac p4$ --- 正方形。

\emph{面积固定。} 边长为 $x$ 与 $\frac Ax$ 的矩形，周长为
$P(x) = 2\left(x + \frac Ax\right)$，且
$P'(x) = 2\left(1 - \frac{A}{x^2}\right)$，当 $x < \sqrt A$ 时为负，
之后为正，故最小值在 $x = \sqrt{A}$ 处取得，仍是正方形。

\emph{关系。} 两个结论互为对偶。假设一个非正方形矩形在
面积 $A$ 固定时周长最小；由第一个结论，与其周长相同的正方形
面积 $> A$，将其位似缩小至面积 $A$ 会
严格减小其周长，与最小性矛盾。
因此两个结论互相蕴含。
\end{solution}

\begin{solution}{pb:g12:deriv:1}
\textbf{1.} $30x\left(3x^2 + 1\right)^4$; \quad
$\dfrac{x}{\sqrt{x^2 + 9}}$; \quad
$-\dfrac{2x}{\left(x^2 + 1\right)^2}$.

\textbf{2.} $y(4) = 5$，斜率 $\frac45$，切线为
$y = 5 + \frac45(x - 4)$，即 $y = \frac45 x + \frac95$。

\textbf{3.} $f'(x) = \frac{(x^2 + 1) - x \cdot 2x}{(x^2+1)^2}
= \frac{1 - x^2}{(x^2 + 1)^2}$：符号依次为负、正、负，
以 $-1$ 和 $1$ 为分界：最小值 $f(-1) = -\frac12$，最大值
$f(1) = \frac12$，且两侧水平渐近线均为 $0$。

\textbf{4.} $g'(x) = 3x^2 - 6x = 3x(x - 2)$：递增、
递减、递增，以 $0$ 和 $2$ 为分界；极大值
$g(0) = 4$，极小值 $g(2) = 0$。$g''(x) = 6x - 6$：在
$1$ 之前为凹，之后为凸，拐点为 $(1, 2)$。

\textbf{5.} 在 $1$ 处的切线，斜率 $g'(1) = -3$：
$y = 2 - 3(x - 1) = -3x + 5$。作差：
$x^3 - 3x^2 + 4 - (-3x + 5) = x^3 - 3x^2 + 3x - 1 =
(x - 1)^3$，在 $1$ 处变号，切线
\emph{穿过}曲线 --- 拐点处特有的
行为，曲线在此切换到切线的另一侧。

\textbf{6.} $T(x) = \dfrac{\sqrt{x^2 + 900}}{5} +
\dfrac{\sqrt{(40 - x)^2 + 400}}{2}$。在 $0$ 之前入水或
超过 $40$，都会同时加长\emph{两段}路程：毫无意义。

\textbf{7.} $T'(x) = \dfrac{x}{5\sqrt{x^2 + 900}} -
\dfrac{40 - x}{2\sqrt{(40 - x)^2 + 400}}$.

\textbf{8.} 在跑步段构成的三角形中，沿岸的边
长为 $x$，斜边为 $\sqrt{x^2 + 900}$：与垂线夹角的正弦
恰为两者之比
（对边比斜边）；同理
$\sin\theta_2 = \frac{40 - x}{\sqrt{(40-x)^2 + 400}}$。于是
$T'(x) = 0$ 即
$\frac{\sin\theta_1}{5} = \frac{\sin\theta_2}{2}$ --- 斯涅尔
定律，速度分别为 $5$ 与 $2$。

\textbf{9.} $T(24) \approx 20.5$ 秒；$T(40) = 10 + 10 = 20$ 秒；
$T(0) = 6 + \sqrt{2000}/2 \approx 28.4$ 秒。直线路径
仍输于“先全程跑、再直游”的方案 --- 两者都输给
折射路径。

\textbf{10.} 对 $T'$ 用二分法（在 $24$ 处为负，在
$40$ 处为正）收敛到 $x^* \approx 34$ 米，
$T(x^*) \approx 19.5$ 秒，比直线路径快约一秒 ---
救援与悲剧之间的
差别。

\textbf{11.} 每段用时都是 $x$ 的凸函数
（$\sqrt{\text{二次式}}$ 型分支），因此 $T$ 为凸函数，而
可导凸函数的临界点即其全局最小值
（见\cref{exo:g12:deriv:8}）：$x^*$ 不仅是
驻点，更是无可匹敌的。

\textbf{12.} 光在水中传播较慢；由费马原理，
从空气到水的最快路径恰在水面处按
斯涅尔定律折射 --- 因此水下铅笔发出的光线到达眼睛时
已经弯折，而大脑依直线外推，看到铅笔
在水面处像是被折断了。

\textbf{13.} $(x^2)'' = 2 > 0$：凸。由
中点处的凸性，即 $f\left(\frac{a+b}{2}\right) \leq
\frac{f(a) + f(b)}{2}$，正是所要证明的。直接验算：
$\frac{a^2 + b^2}{2} - \left(\frac{a+b}{2}\right)^2 =
\frac{(a - b)^2}{4} \geq 0$.

\textbf{14.} 对第 13 题两边取平方根（递增）：
$\frac{a + b}{2} \leq \sqrt{\frac{a^2 + b^2}{2}}$。取 $2, 8$，完整的不等链为
调和平均 $\frac{2 \cdot 16}{10} = 3.2$；几何平均
$\sqrt{16} = 4$；算术平均 $5$；平方平均
$\sqrt{34} \approx 5.83$：每一个都不大于下一个。

\textbf{15.} $\frac1x$ 在 $1$ 处的切线：值为 $1$，斜率
$-1$：$y = 2 - x$。凸曲线总位于其切线
上方（见\cref{thm:g12:deriv:convexity}）：$\frac1x \geq 2 - x$ 对所有
$x > 0$ 成立，等号恰在切点 $x = 1$ 处取得。

\textbf{16.} 在拐点处，\emph{每日}新增病例数
（即导数）达到峰值：增长不再加速，开始
减速 --- 这是第一个数学信号，表明疫情
正在转向，远早于病例数本身下降。拐点之前，
$f'' > 0$（每天都比前一天更严重）；拐点之后，
$f'' < 0$（仍在增长，但增速放缓）。

\textbf{17.} 由 $S'(r) = 4\pi r - \frac{2V}{r^2} = 0$ 得
$r^3 = \frac{V}{2\pi} = \frac{330}{2\pi}$:
$r \approx 3.74$ 厘米，而
$h = \frac{330}{\pi r^2} \approx 7.49$ 厘米。

\textbf{18.} $h = \frac{V}{\pi r^2} = \frac{2\pi r^3}{\pi r^2}
= 2r$：最优时高度总等于直径，无论
体积多大。

\textbf{19.} $S''(r) = 4\pi + \frac{4V}{r^3} > 0$：凸，故
临界半径即为全局最小值。现实中的罐头做得更高，
是为了便于握持、堆叠、货架展示，以及顶部与底部
的金属更厚 --- 优化总是精确地最小化你所写下的，
而非你的本意。

\textbf{20.} 救生员问题：建立 $T(x)$ 的模型，求导，运用斯涅尔
定律，判断凸性，节省一秒。罐头问题：建立 $S(r)$ 的模型，求导，
$h = 2r$，判断凸性，节省金属。梁问题
（见\cref{pb:g11:deriv:1}）：建立 $S(w)$ 的模型，求导，
$d = w\sqrt2$，列表，刚度有所增加。再由费马
原理，光本身在它所遇到的每一个界面上
都在执行这一流程 --- 大自然才是最早的优化者。
\end{solution}
